% Security basics (reuse tutorial slides), mapped to automotive
\xchapter{Security, in a Car}

\begin{frame}{Safety vs.\ Security: Two Different Enemies}
  \begin{columns}[T]
    \begin{column}{0.5\textwidth}
      \textbf{Safety} \;(ISO 26262)
      \begin{itemize}
        \item Guards against \emph{bad luck}: random faults, design errors.
        \item Assumes nobody is out to get you.
        \item Mature in cars since the 2000s.
      \end{itemize}
    \end{column}
    \begin{column}{0.5\textwidth}
      \textbf{Security} \;(ISO/SAE 21434)
      \begin{itemize}
        \item Guards against \emph{bad intent}: a thinking adversary.
        \item Assumes someone hunts the weakest link.
        \item Only mandatory since 2022 (UN R155).
      \end{itemize}
    \end{column}
  \end{columns}
  \vspace{0.3cm}
  \begin{alertblock}{The key insight}
    A car can be perfectly \emph{safe}, failing predictably, and still wide open, because the attacker picks when and how it fails. Modern cars have to do both.
  \end{alertblock}
\end{frame}
